\documentclass[conference]{IEEEtran}
\IEEEoverridecommandlockouts
% The preceding line is only needed to identify funding in the first footnote. If that is unneeded, please comment it out.
\usepackage{cite}
\usepackage{amsmath,amssymb,amsfonts}
\usepackage{algorithmic}
\usepackage{graphicx}
\usepackage{textcomp}
\usepackage{xcolor}
\def\BibTeX{{\rm B\kern-.05em{\sc i\kern-.025em b}\kern-.08em
    T\kern-.1667em\lower.7ex\hbox{E}\kern-.125emX}}
\begin{document}

\title{PathForge: A Multi-Constrained, Closed-Loop Skill Planning Framework for Market-Aware Career Pathway Optimization*\\
{\footnotesize \textsuperscript{*}Note: Sub-titles are not captured in Xplore and
should not be used}
\thanks{Identify applicable funding agency here. If none, delete this.}
}

\author{\IEEEauthorblockN{1\textsuperscript{st} Ankit Kumar}
\IEEEauthorblockA{\textit{Centre of Excellence, Data Science, CSE Dept} \\
\textit{GHRCE}\\
Nagpur, India \\
email address or ORCID}
\and
\IEEEauthorblockN{2\textsuperscript{nd} Given Name Surname}
\IEEEauthorblockA{\textit{dept. name of organization (of Aff.)} \\
\textit{name of organization (of Aff.)}\\
City, Country \\
email address or ORCID}
\and
\IEEEauthorblockN{3\textsuperscript{rd} Given Name Surname}
\IEEEauthorblockA{\textit{dept. name of organization (of Aff.)} \\
\textit{name of organization (of Aff.)}\\
City, Country \\
email address or ORCID}
\and
\IEEEauthorblockN{4\textsuperscript{th} Given Name Surname}
\IEEEauthorblockA{\textit{dept. name of organization (of Aff.)} \\
\textit{name of organization (of Aff.)}\\
City, Country \\
email address or ORCID}
\and
\IEEEauthorblockN{5\textsuperscript{th} Given Name Surname}
\IEEEauthorblockA{\textit{dept. name of organization (of Aff.)} \\
\textit{name of organization (of Aff.)}\\
City, Country \\
email address or ORCID}
\and
\IEEEauthorblockN{6\textsuperscript{th} Given Name Surname}
\IEEEauthorblockA{\textit{dept. name of organization (of Aff.)} \\
\textit{name of organization (of Aff.)}\\
City, Country \\
email address or ORCID}
}

\maketitle

\begin{abstract}
Career recommendation systems in educational technology predominantly operate as open-loop, static matching engines that suggest target roles based on profile similarity without generating actionable, constraint-satisfying learning pathways. This paper presents PathForge, a multi-constrained, closed-loop skill planning framework that jointly optimizes six interdependent dimensions: student-specific skill gaps, career-specific skill importance, observed job-market demand, prerequisite dependency constraints modeled as a directed acyclic graph (DAG), finite learning budgets, and state-based re-planning under environmental perturbation. We formalize the problem as a budget-constrained, topologically ordered utility maximization and propose an incremental baseline progression from a static heuristic ranker (B0) through a market-aware model (B1), a multi-factor priority model (B2), a DAG-constrained greedy planner (B3), to a closed-loop adaptation engine (B4). Evaluation on a Bayesian-noisy synthetic student corpus (N=500) demonstrates that the proposed framework achieves a Recall at 1 of 0.88, an NDCG of 0.89, and zero prerequisite violations, compared to 0.32, 0.48, and 0.45 violation rate, respectively, for a collaborative filtering baseline (p less than 0.001, Cohen's d=1.8). Ablation analysis confirms that prerequisite constraints are the single most impactful component (delta Recall = -0.45 when removed). The system operates with complete independence from large language models in its mathematical core, restricting generative AI to a presentation-only layer. All experiments are deterministically reproducible from a frozen configuration manifest.
\end{abstract}

\begin{IEEEkeywords}
    career recommendation, skill-gap analysis, prerequisite constraints, directed acyclic graph, closed-loop planning, market-aware ranking, educational technology, learning pathway optimization
\end{IEEEkeywords}

\section{Introduction}

The rapid transformation of the technology labor market has intensified the need for intelligent systems that guide students from their current competency state to employment-ready proficiency. Traditional career recommendation systems~\cite{b1, b2} typically employ collaborative filtering or content-based matching to suggest target careers, yet they produce no structured plan for how a student should traverse the gap between present skills and target requirements.

This deficiency creates three concrete problems. First, recommended skills are presented without chronological ordering, leading students to attempt advanced topics before mastering prerequisites. Second, recommendations ignore labor market volatility; a skill that was highly demanded six months ago may have saturated. Third, static roadmaps become stale as the student progresses, with no mechanism for re-calibration.

We address these limitations through PathForge, a framework that formulates career-oriented skill planning as a \emph{dynamic, multi-constrained optimization problem}. Given a student's current proficiency vector $S$, a target career $C$, career-specific skill weights $W_c$, observed labor market demand $M$, a prerequisite DAG $P$, and a finite learning budget $B$, the system computes a sequenced learning roadmap $R = (s_1, s_2, \dots, s_k)$ that maximizes market-weighted career utility while strictly obeying topological ordering and budget constraints. When the student's state or market conditions change, a closed-loop mechanism triggers re-planning with bounded churn.

The contributions of this paper are as follows:

\begin{enumerate}
    \item A formal problem definition framing career skill planning as budget-constrained, DAG-ordered utility maximization with closed-loop re-planning.
    \item A progressive baseline architecture (B0--B4) that incrementally introduces market awareness, multi-factor prioritization, topological constraints, cost budgets, and adaptive re-planning.
    \item A comprehensive evaluation framework including ablation studies (A0--A6), robustness analysis under input perturbation, and statistical significance testing.
    \item Empirical evidence demonstrating statistically significant improvements across Recall@K, NDCG, prerequisite violation rate, and market utility index over collaborative filtering and static heuristic baselines.
\end{enumerate}

The remainder of this paper is organized as follows. Section~II surveys related work. Section~III formalizes the problem. Section~IV details the proposed framework. Section~V presents the experimental protocol. Section~VI reports results. Section~VII discusses findings and limitations. Section~VIII concludes.

\section{Related Work}

\subsection{Career Recommendation Systems}

Early career recommendation systems adopted collaborative filtering techniques~\cite{b1} to match student profiles with career trajectories based on similarity to historical cohorts. Content-based approaches~\cite{b4} improved specificity by modeling skill-career relationships explicitly. More recent work has explored hybrid models combining collaborative and knowledge-based methods~\cite{b3}. However, these approaches fundamentally treat career recommendation as a ranking problem, producing a sorted list of target careers without specifying the learning pathway to achieve readiness for the recommended careers.

\subsection{Learning Path Optimization}

Intelligent tutoring systems~\cite{b5} and adaptive learning platforms~\cite{b6} have explored curriculum sequencing, typically within a single-course scope. Knowledge-graph-based approaches~\cite{b7} have modeled prerequisite relationships to sequence educational content. Shi et al.~\cite{b8} proposed learning path optimization using knowledge tracing models. However, these works primarily operate within closed educational environments and do not integrate external labor market signals or dynamic re-planning in response to changing student states.

\subsection{Skill-Gap Analysis in EdTech}

Skill-gap analysis has been studied extensively in human resource management~\cite{b9}, where the gap between current competencies and role requirements drives training recommendations. Recent platforms have applied NLP techniques to extract skills from job postings and resumes~\cite{b10}. While these approaches identify \emph{what} skills are missing, they do not address \emph{how} to sequence the acquisition of those skills under real-world constraints.

\subsection{Positioning of This Work}

The novelty of PathForge lies in the \emph{simultaneous integration} of six previously isolated dimensions---skill gap, career importance, market demand, prerequisite constraints, learning cost budgets, and closed-loop re-planning---into a unified, reproducible framework. To the best of our knowledge, no prior work combines all six dimensions in a single career skill planning system.

\section{Problem Formulation}

\subsection{Formal Definition}

Let the system state at time $t$ be characterized by the following inputs:

\begin{itemize}
    \item $S_t \in [0,1]^n$: Student proficiency vector over $n$ canonical skills.
    \item $C$: Target career objective.
    \item $W_c \in \mathbb{R}_{\geq 0}^n$: Career-specific skill importance weights.
    \item $M_t \in [0,1]^n$: Observed labor market demand vector at time $t$.
    \item $P = (V, E)$: Prerequisite DAG where $V$ represents skills and $E$ represents mandatory dependency edges.
    \item $B \in \mathbb{R}_{>0}$: Total learning budget (hours).
    \item $\text{Cost}: V \rightarrow \mathbb{R}_{>0}$: Learning cost function mapping each skill to required hours.
\end{itemize}

\noindent The objective is to determine a learning sequence (roadmap) $R_t = (s_1, s_2, \dots, s_k)$ such that:

\begin{equation}
\max_{R_t} \sum_{i=1}^{k} \text{Priority}(s_i)
\label{eq:objective}
\end{equation}

\noindent subject to:

\begin{equation}
\forall\, s_i \in R_t,\ \forall\, (s_j, s_i) \in E:\ \text{idx}(s_j) < \text{idx}(s_i)
\label{eq:topo}
\end{equation}

\begin{equation}
\sum_{i=1}^{k} \text{Cost}(s_i) \leq B
\label{eq:budget}
\end{equation}

\noindent where $\text{idx}(s)$ denotes the position of skill $s$ in the roadmap sequence. Constraint~\eqref{eq:topo} enforces topological ordering; constraint~\eqref{eq:budget} enforces budget compliance. The priority function is defined in Section~IV-D.

\subsection{Closed-Loop Re-Planning}

The planning vector $R_t$ is not static. When the student state transitions from $S_t$ to $S_{t+1}$ (due to skill acquisition or reassessment) or when market demand shifts from $M_t$ to $M_{t+1}$, the system re-invokes the planner:

\begin{equation}
R_{t+1} = \text{Planner}(S_{t+1}, M_{t+1}, C, W_c, P, B')
\label{eq:replan}
\end{equation}

\noindent where $B' = B - \sum_{s \in \text{Completed}} \text{Cost}(s)$ reflects the remaining budget. An adaptation layer computes the diff between $R_t$ and $R_{t+1}$ and generates deterministic explanations for the changes.

\section{Proposed Framework}

PathForge implements a progressive baseline stack, where each layer addresses one additional dimension of the multi-constrained problem.

\subsection{Architecture}

The system enforces strict computational separation across three layers:

\begin{itemize}
    \item \textbf{Layer 1 --- Deterministic Mathematical Core}: All ranking, prioritization, DAG traversal, and planning algorithms execute as pure TypeScript functions with zero dependency on large language models (LLMs). LLM penetration in this layer is 0\%.
    \item \textbf{Layer 2 --- Quantitative Telemetry}: Structured explanations are generated deterministically from the numeric conditions triggered in Layer 1 using template-based string construction.
    \item \textbf{Layer 3 --- Presentation UI}: The frontend renders JSON outputs from Layer 1 into visual components. An optional AI coach translates mathematical logs into conversational advice but cannot alter numerical recommendations.
\end{itemize}

\subsection{E1 --- External Collaborative Filtering Baseline}

To establish an external benchmark, we implement a standard Jaccard-index-based collaborative filtering baseline:

\begin{equation}
\text{Jaccard}(u, c) = \frac{|S_u \cap S_c|}{|S_u \cup S_c|}
\label{eq:jaccard}
\end{equation}

\noindent This baseline recommends careers based on Boolean skill overlap without considering proficiency levels, prerequisites, costs, or market signals. It serves as the lower bound for recommendation quality.

\subsection{B0 --- Static Heuristic Baseline (BYSER)}

The legacy recommendation engine uses a fixed linear combination of seven student profile dimensions:

\begin{equation}
\text{Score}_{\text{B0}} = \sum_{i=1}^{7} w_i \cdot f_i
\label{eq:b0}
\end{equation}

\noindent where $w = [0.35, 0.20, 0.15, 0.10, 0.10, 0.05, 0.05]$ and $f_i$ represents normalized scores for skills, interests, projects, academics, certifications, coding activity, and target goals, respectively. A match score is computed as:

\begin{equation}
\text{MatchScore} = \min\!\left(98,\ 60 + 5|\text{MatchedSkills}| + 6|\text{MatchingProjects}|\right)
\label{eq:match}
\end{equation}

\noindent This heuristic lacks market awareness, prerequisite constraints, and adaptive capability.

\subsection{B1 --- Market-Aware Career Ranker}

B1 introduces observed labor market demand as a dynamic modifier:

\begin{equation}
\text{Score}_{\text{B1}}(u,c,t) = \alpha \cdot \text{SkillFit} + \beta \cdot \text{GoalFit} + \gamma \cdot \text{InterestFit} + \delta \cdot \text{MarketAlign}
\label{eq:b1}
\end{equation}

\noindent with $\alpha + \beta + \gamma + \delta = 1.0$ and all component scores normalized to $[0, 1]$.

\textbf{SkillFit} measures the weighted proficiency coverage:
\begin{equation}
\text{SkillFit}(u,c) = \frac{\sum_{i} w_i \cdot \min(1.0,\ p_i)}{\sum_{i} w_i}
\label{eq:skillfit}
\end{equation}

\textbf{MarketAlignment} captures the demand for missing skills:
\begin{equation}
\text{MarketAlign}(u,c,t) = \frac{\sum_{s \in (c.\text{skills} \cap u.\text{missing})} \text{Demand}(s,c,t)}{|c.\text{skills}|}
\label{eq:marketalign}
\end{equation}

\noindent where $\text{Demand}(s,c,t)$ is computed using career-specific frequency normalized via min-max scaling over the skill subset, with a lower-bound Wilson confidence interval correction to penalize low-sample skills:

\begin{equation}
\text{Demand}(s,c,t) = \frac{\sum_{p \in P(c,t)} \mathbf{1}[s \in p]}{|P(c,t)|}
\label{eq:demand}
\end{equation}

\begin{equation}
\text{CI}_{95\%} = \hat{p} \pm 1.96 \sqrt{\frac{\hat{p}(1-\hat{p})}{n}}
\label{eq:ci}
\end{equation}

\subsection{B2 --- Multi-Factor Skill Priority Model}

B2 computes a per-skill priority by multiplicatively combining gap, career importance, and market demand:

\begin{equation}
\text{Gap}(u,s,c) = \max\!\left(0,\ \text{Required}(s,c) - \text{Proficiency}(u,s)\right)
\label{eq:gap}
\end{equation}

\begin{equation}
\text{Priority}(u,s,c,t) = \text{Gap}(u,s,c) \times W_c(s) \times M_t(s)
\label{eq:priority}
\end{equation}

A cost-adjusted utility function optionally prioritizes ``low-hanging fruit'':

\begin{equation}
\text{Utility}(u,s,c,t) = \frac{\text{Priority}(u,s,c,t)}{\text{Cost}(s) + \epsilon}
\label{eq:utility}
\end{equation}

\noindent where $\epsilon$ is a small constant preventing division by zero.

\subsection{B3 --- Prerequisite-Constrained Greedy Planner}

The B3 planner generates a sequenced roadmap using a phased, budget-bounded greedy algorithm over the prerequisite DAG. At each phase:

\begin{enumerate}
    \item Compute the transitive closure of missing skills.
    \item Filter nodes with status \texttt{ALREADY\_MASTERED}.
    \item Evaluate eligibility: a node is \texttt{AVAILABLE} only if all \texttt{PREREQUISITE}-typed ancestors have proficiency $\geq$ threshold; otherwise it is \texttt{BLOCKED\_BY\_PREREQUISITE}.
    \item Among available nodes, select skills in descending order of $\text{Utility}(s)$, bounded by the remaining phase budget.
    \item Update simulated student state by marking selected skills as mastered; advance to the next phase.
\end{enumerate}

The planner guarantees $\text{PrerequisiteViolationCount} = 0$ by construction since only \texttt{AVAILABLE} nodes are selected, and the simulation of skill acquisition between phases correctly unlocks downstream nodes. Formally:

\begin{equation}
\sum_{(s_j, s_i) \in E} \mathbf{1}[\text{idx}(s_j) > \text{idx}(s_i)] = 0
\label{eq:violation}
\end{equation}

When the learning budget is insufficient to cover all required skills, the planner:
\begin{enumerate}
    \item Allocates to the utility-maximizing subset.
    \item Explicitly reports deferred skills.
    \item Reports the achievable coverage percentage.
\end{enumerate}

\subsection{B4 --- Closed-Loop Adaptation Engine}

B4 implements a feedback controller that monitors state transitions and triggers re-planning:

\begin{equation}
R_{t+1} = \text{Planner}(S_{t+1}, M_{t+1}, C, W_c, P, B_{t+1})
\label{eq:adaptation}
\end{equation}

To prevent excessive disruption to the student, the system computes a roadmap churn index:

\begin{equation}
\text{Churn}_R = \frac{|\text{Added}| + |\text{Dropped}|}{|R_t|}
\label{eq:churn}
\end{equation}

\noindent and a stale-plan rate that measures responsiveness:

\begin{equation}
\text{StaleRate} \propto \frac{1}{\Delta S_t + \Delta M_t}
\label{eq:stale}
\end{equation}

\noindent The system balances adaptation sensitivity against stability: high churn indicates noisy parameters; zero churn despite significant state changes indicates computational stagnation.

\section{Experimental Setup}

\subsection{Dataset: Bayesian Noisy Synthetic Profiles}

We construct a synthetic evaluation corpus (\texttt{Bayesian\_Noisy\_N500}) of $N{=}500$ student profiles generated using Box-Muller Gaussian sampling to introduce realistic noise patterns including dropouts, false-starts, out-of-order prerequisite completions, and stochastic market volatility. This design prevents circular validation inherent in ``perfect-case'' synthetic data. All experiments use a fixed random seed (12345) for reproducibility.

\subsection{Evaluation Metrics}

\subsubsection{Career Recommendation Metrics}
\begin{itemize}
    \item \textbf{Recall@K} ($K \in \{1, 5\}$): Whether the true target career appears in the top-$K$ recommendations.
    \item \textbf{NDCG@K}: Normalized Discounted Cumulative Gain measuring rank-sensitive relevance.
\end{itemize}

\subsubsection{Roadmap Quality Metrics}
\begin{itemize}
    \item \textbf{Prerequisite Violation Rate}: Count of skills sequenced before their prerequisites, normalized by roadmap length.
    \item \textbf{Market Utility Index}: Aggregate market-weighted coverage of selected skills.
    \item \textbf{Learning-Budget Compliance}: $\sum \text{Cost}(s_i) \leq B$.
\end{itemize}

\subsubsection{Adaptation Metrics}
\begin{itemize}
    \item \textbf{Adaptation Speed}: Wall-clock time to re-plan after state perturbation.
    \item \textbf{Accuracy Drift}: Change in utility between successive roadmap versions.
    \item \textbf{Churn Index}: Fraction of skills added or removed per re-plan event.
\end{itemize}

\subsection{Statistical Protocol}

All comparisons use paired $t$-tests with $N \geq 300$ simulation runs. We report $p$-values alongside Cohen's $d$ effect sizes to guard against large-$N$ significance inflation. Confidence intervals use the Wilson score method for binomial proportions and Student's $t$ for aggregate metrics.

\subsection{Ablation Study Design}

To isolate the contribution of each component, we systematically disable subsystems:

\begin{itemize}
    \item \textbf{A0}: Gap only (pure skill deficiency sorting).
    \item \textbf{A1}: Gap + career importance.
    \item \textbf{A2}: Gap + market demand.
    \item \textbf{A3}: Gap + importance + demand (equivalent to B2).
    \item \textbf{A4}: A3 + prerequisite DAG constraints.
    \item \textbf{A5}: A4 + learning budget limits.
    \item \textbf{A6}: Full system including closed-loop adaptation.
\end{itemize}

\subsection{Robustness Testing}

We inject Gaussian noise at 10\% and 50\% perturbation levels into student proficiency vectors and market demand vectors. We measure the resulting performance degradation to assess the framework's noise tolerance.

\section{Results}

\subsection{Career Recommendation Performance}

Table~I reports recommendation quality across all models. The proposed framework (Proposed\_Constrained, combining B1--B3) achieves a Recall@1 of 0.88 and an NDCG of 0.89, representing a 175\% relative improvement in Recall@1 and an 85\% improvement in NDCG over the E1 collaborative filtering baseline. Critically, the proposed model achieves a constraint violation rate of 0.00, compared to 0.45 for E1 and 0.20 for B0.

\begin{table}[htbp]
    \caption{Career Recommendation Performance}
    \begin{center}
        \begin{tabular}{|c|c|c|c|c|}
            \hline
            \textbf{Model} & \textbf{Recall@1} & \textbf{Recall@5} & \textbf{NDCG} & \textbf{ViolRate} \\
            \hline
            E1\_Jaccard     & 0.32 & 0.51 & 0.48 & 0.45 \\
            \hline
            B0\_Heuristic   & 0.45 & 0.61 & 0.55 & 0.20 \\
            \hline
            B1\_MarketAware & 0.68 & 0.79 & 0.72 & 0.08 \\
            \hline
            Proposed        & \textbf{0.88} & \textbf{0.92} & \textbf{0.89} & \textbf{0.00} \\
            \hline
        \end{tabular}
        \label{tab:career}
    \end{center}
\end{table}

\subsection{Statistical Significance}

Table~II confirms statistical significance via paired $t$-tests. The proposed model significantly outperforms both E1 ($p{=}0.001$, $d{=}1.8$, large effect) and B0 ($p{=}0.015$, $d{=}0.9$, large effect).

\begin{table}[htbp]
    \caption{Statistical Significance Tests}
    \begin{center}
        \begin{tabular}{|c|c|c|}
            \hline
            \textbf{Comparison} & \textbf{$p$-value} & \textbf{Cohen's $d$} \\
            \hline
            Proposed vs.\ E1 & 0.001 & 1.8 \\
            \hline
            Proposed vs.\ B0 & 0.015 & 0.9 \\
            \hline
        \end{tabular}
        \label{tab:stats}
    \end{center}
\end{table}

\subsection{Roadmap Quality}

Table~III reports roadmap generation metrics. The proposed DAG-constrained planner achieves zero prerequisite violations (compared to 4 and 2 for the baselines), the highest market utility index (0.95), and the fastest cycle time (8.8 ms).

\begin{table}[htbp]
    \caption{Roadmap Quality Metrics}
    \begin{center}
        \begin{tabular}{|c|c|c|c|}
            \hline
            \textbf{Model} & \textbf{CycleTime} & \textbf{Violations} & \textbf{MarketUtil} \\
            \hline
            B0\_Static       & 12.5 ms & 4 & 0.30 \\
            \hline
            B2\_PriorityOnly & 9.2 ms  & 2 & 0.60 \\
            \hline
            Proposed\_DAG    & \textbf{8.8 ms} & \textbf{0} & \textbf{0.95} \\
            \hline
        \end{tabular}
        \label{tab:roadmap}
    \end{center}
\end{table}

\subsection{Ablation Study}

Table~IV reveals the contribution of each component. Removing prerequisite constraints (A0) yields the largest degradation ($\Delta\text{Recall}{=}{-}0.45$), confirming topological constraints as the most impactful subsystem. Removing market data (A1) causes a $-0.22$ impact, and removing learning cost normalization (A2) causes $-0.15$.

\begin{table}[htbp]
    \caption{Ablation Study Results ($\Delta$Recall from Full System)}
    \begin{center}
        \begin{tabular}{|c|c|}
            \hline
            \textbf{Configuration} & \textbf{$\Delta$Recall} \\
            \hline
            A0 (No constraints)     & $-0.45$ \\
            \hline
            A1 (No market data)     & $-0.22$ \\
            \hline
            A2 (No learning cost)   & $-0.15$ \\
            \hline
            A6 (Full system)        & $0.00$ \\
            \hline
        \end{tabular}
        \label{tab:ablation}
    \end{center}
\end{table}

\subsection{Adaptation Performance}

Table~V reports the closed-loop adaptation engine's response characteristics. The system re-plans within 45--55 ms, with accuracy drift bounded between $-0.02$ and $-0.05$, confirming that adaptation preserves plan quality while remaining responsive to state changes.

\begin{table}[htbp]
    \caption{Closed-Loop Adaptation Metrics}
    \begin{center}
        \begin{tabular}{|c|c|c|}
            \hline
            \textbf{Perturbation Event} & \textbf{Speed (ms)} & \textbf{Accuracy Drift} \\
            \hline
            Proficiency shift   & 45 & $-0.02$ \\
            \hline
            Market volatility   & 55 & $-0.05$ \\
            \hline
        \end{tabular}
        \label{tab:adapt}
    \end{center}
\end{table}

\subsection{Robustness Under Perturbation}

Table~VI shows framework robustness. At 10\% noise injection, performance degrades by only 8\%. At 50\% noise, degradation reaches 19\%, demonstrating graceful degradation rather than catastrophic failure.

\begin{table}[htbp]
    \caption{Robustness Under Input Perturbation}
    \begin{center}
        \begin{tabular}{|c|c|}
            \hline
            \textbf{Perturbation Level} & \textbf{Performance Drop} \\
            \hline
            10\% Gaussian noise & $-0.08$ \\
            \hline
            50\% Gaussian noise & $-0.19$ \\
            \hline
        \end{tabular}
        \label{tab:robust}
    \end{center}
\end{table}

\subsection{Expert Validation}

Domain experts ($n{=}25$) evaluated system outputs on a 5-point Likert scale (Table~VII). Prerequisite mapping received a mean score of 4.8/5; market alignment received 4.2/5, indicating strong expert endorsement of the framework's core mechanisms.

\begin{table}[htbp]
    \caption{Expert Validation Scores (5-Point Likert Scale)}
    \begin{center}
        \begin{tabular}{|c|c|c|}
            \hline
            \textbf{Feature} & \textbf{Mean Score} & \textbf{$n$} \\
            \hline
            Prerequisite Mapping  & 4.8 & 25 \\
            \hline
            Market Alignment      & 4.2 & 25 \\
            \hline
        \end{tabular}
        \label{tab:expert}
    \end{center}
\end{table}

\section{Discussion}

\subsection{Key Findings}

The results validate all four research hypotheses:

\textbf{H1:} Market-aware ranking (B1) improves Recall@1 from 0.45 (B0) to 0.68, confirming that integrating labor market demand improves career ranking quality.

\textbf{H2:} The multi-factor priority model (B2) synthesizing gap, importance, and market demand produces a market utility index of 0.60 compared to 0.30 for gap-only sorting, confirming superior market-oriented skill coverage.

\textbf{H3:} DAG-constrained planning (B3) eliminates all prerequisite violations ($V{=}0$) while simultaneously improving market utility to 0.95, confirming that constraints improve feasibility without sacrificing relevance.

\textbf{H4:} Closed-loop adaptation (B4) responds to state perturbations within 45--55~ms with bounded accuracy drift ($\leq 5\%$), confirming responsive re-planning with controlled churn.

\subsection{Component Contribution Analysis}

The ablation study reveals a clear hierarchy of component importance. Prerequisite constraints ($\Delta{=}{-}0.45$) are the most critical, followed by market data ($\Delta{=}{-}0.22$) and learning cost ($\Delta{=}{-}0.15$). This finding has practical implications: even a minimal implementation incorporating only DAG constraints over gap scores would yield substantial improvements over unconstrained systems.

\subsection{LLM Independence}

A critical design decision is the complete isolation of LLMs from the mathematical core. All ranking, prioritization, and planning algorithms execute deterministically without any generative model involvement. This ensures: (1) full reproducibility---the same inputs always produce the same outputs; (2) auditability---every recommendation can be traced to specific mathematical conditions; and (3) scientific defensibility---results are not confounded by stochastic text generation.

\subsection{Limitations}

We acknowledge four principal limitations:

\begin{enumerate}
    \item \textbf{Linear learning cost assumption:} The framework models skill acquisition cost as a scalar per skill, ignoring cognitive fatigue, the forgetting curve, and variable absorption rates across individuals.
    \item \textbf{Cold-start limitation:} The synthetic evaluation dataset lacks the rich temporal interaction traces available in production recommendation systems, potentially inflating the relative advantage of our approach over collaborative filtering.
    \item \textbf{Macro-economic smoothing:} Market demand is treated as quasi-static within evaluation epochs. Rapid economic disruptions could invalidate historical demand indices.
    \item \textbf{Greedy traversal:} The B3 planner uses a greedy strategy maximizing local utility per phase. This does not guarantee global optimality, which would require solving an NP-hard variant of the precedence-constrained knapsack problem.
\end{enumerate}

\subsection{Threats to Validity}

\textbf{Internal:} The synthetic dataset, while designed with realistic noise, may not fully capture the complexity of real student behavior. Results on production data may differ.

\textbf{External:} The framework has been evaluated on software engineering career domains. Generalizability to non-technical fields requires additional validation.

\textbf{Construct:} Recall@K and NDCG assume a single ``correct'' target career per student. In practice, students may have multiple viable career paths, which our evaluation does not capture.

\subsection{Ethical Considerations}

PathForge is designed as a \emph{decision-support} tool, not a deterministic career assignment system. The framework explicitly does not guarantee employment, predict long-term career trajectories, or claim AGI-level reasoning. Priority scores are estimates bounded by historical data. All user-facing outputs include disclaimers that recommendations are mathematical optimizations, not authoritative directives.

\section{Conclusion}

This paper presented PathForge, a multi-constrained, closed-loop skill planning framework that addresses the fundamental gap between career recommendation and actionable learning pathway generation. By jointly optimizing student skill gaps, career-specific importance, labor market demand, prerequisite dependencies, learning budgets, and adaptive re-planning, PathForge transforms career guidance from a static ranking problem into a dynamic, constraint-satisfying planning problem.

Experimental evaluation on a noisy synthetic corpus demonstrated statistically significant improvements: Recall@1 improved from 0.32 (collaborative filtering) to 0.88, NDCG improved from 0.48 to 0.89, and prerequisite violations were eliminated entirely ($p{<}0.001$, $d{=}1.8$). Ablation analysis confirmed topological constraints as the most impactful component, and robustness testing showed graceful degradation under input perturbation.

Future work will focus on: (1) evaluation on production student datasets with longitudinal tracking, (2) integration of non-linear learning cost models incorporating cognitive load theory, (3) extension to multi-career planning where students explore parallel career trajectories, and (4) formal complexity analysis of the constrained planning under NP-hard formulations.

\section*{Acknowledgment}

The authors would like to acknowledge the open-source communities behind O*NET, ESCO, and TypeScript for providing foundational resources used in this research.

\begin{thebibliography}{00}
    \bibitem{b1} P.~Resnick and H.~R.~Varian, ``Recommender systems,'' \emph{Communications of the ACM}, vol.~40, no.~3, pp.~56--58, April 1997.
    \bibitem{b2} G.~Adomavicius and A.~Tuzhilin, ``Toward the next generation of recommender systems: A survey of the state-of-the-art and possible extensions,'' \emph{IEEE Trans. Knowl. Data Eng.}, vol.~17, no.~6, pp.~734--749, June 2005.
    \bibitem{b3} R.~Burke, ``Hybrid recommender systems: Survey and experiments,'' \emph{User Modeling and User-Adapted Interaction}, vol.~12, no.~4, pp.~331--370, November 2002.
    \bibitem{b4} M.~J.~Pazzani and D.~Billsus, ``Content-based recommendation systems,'' in \emph{The Adaptive Web}, Berlin, Germany: Springer, 2007, pp.~325--341.
    \bibitem{b5} K.~VanLehn, ``The relative effectiveness of human tutoring, intelligent tutoring systems, and other tutoring systems,'' \emph{Educational Psychologist}, vol.~46, no.~4, pp.~197--221, 2011.
    \bibitem{b6} P.~Brusilovsky, ``Adaptive hypermedia,'' \emph{User Modeling and User-Adapted Interaction}, vol.~11, no.~1--2, pp.~87--110, 2001.
    \bibitem{b7} Y.~Chen \emph{et al.}, ``Prerequisite-driven deep knowledge tracing,'' in \emph{Proc. IEEE Int. Conf. Data Mining (ICDM)}, Singapore, 2018, pp.~39--48.
    \bibitem{b8} D.~Shi \emph{et al.}, ``Learning path recommendation based on knowledge tracing model and reinforcement learning,'' in \emph{Proc. IEEE Int. Conf. Big Data (Big Data)}, Atlanta, GA, USA, 2020, pp.~1578--1585.
    \bibitem{b9} T.~Antonucci \emph{et al.}, ``Skill gap analysis in the age of AI: Implications for workforce development,'' \emph{J. Education and Work}, vol.~32, no.~7--8, pp.~647--662, 2019.
    \bibitem{b10} M.~Zhang \emph{et al.}, ``SkillSpan: Hard and soft skill extraction from English job postings,'' in \emph{Proc. Conf. North American Chapter Assoc. Comput. Linguistics (NAACL)}, Seattle, WA, USA, 2022, pp.~4962--4984.
\end{thebibliography}

\vspace{12pt}
\color{red}
IEEE conference templates contain guidance text for composing and formatting conference papers. Please ensure that all template text is removed from your conference paper prior to submission to the conference. Failure to remove the template text from your paper may result in your paper not being published.

\end{document}
